% Lesefassung ohne eigene nummerierte Deklarationen. Alle Aussage-IDs gehören B31.
% symbol-legend.tex wird im Vorspann der Ausgabe geladen.
\NewDocumentCommand\CROrd{s m m}{%
  \mathord{\IfBooleanT{#1}{\SymbolLegendIntroduce{#2}}%
    \SymbolLegendUse{#2}#3}}
\NewDocumentCommand\CRBin{s m m}{%
  \mathbin{\IfBooleanT{#1}{\SymbolLegendIntroduce{#2}}%
    \SymbolLegendUse{#2}#3}}
\NewDocumentCommand\CRRel{s m m}{%
  \mathrel{\IfBooleanT{#1}{\SymbolLegendIntroduce{#2}}%
    \SymbolLegendUse{#2}#3}}
\newcommand\CRK{\CROrd{coordinates-carrier}{\mathcal K_e}}
\newcommand\CRChi{\CROrd{coordinates-map}{\chi_e}}
\newcommand\CRMu{\CROrd{coordinates-multiply}{\mu_e}}
\newcommand\CRLeft{\CROrd{coordinates-left}{Ae}}
\newcommand\CRRight{\CROrd{coordinates-right}{eA}}
\newcommand\CRSquare{\CRBin{coordinates-product}{\square_e}}
\newcommand\CRStar{\CRBin{coordinates-star}{\star}}
\newcommand\CRTimes{\CRBin{coordinates-cartesian}{\times}}
\newcommand\CRComp{\CRBin{coordinates-composition}{\circ}}
\newcommand\CRBij{\CRRel{coordinates-bijection}{\bij}}
\newcommand\CRId[1]{\CROrd{coordinates-identity}{\Id_{#1}}}
\newcommand\CRProof[1]{%
  \href{_B31-coordinates-proofs.pdf\#coordinates.proof.#1}{Beweistabelle}}
\newcommand\CRStatement[3]{%
  \par\Needspace{5\baselineskip}\medskip\noindent
  \textbf{\FormulaRefAuto{#1}[#2]: #3.}\par\smallskip\noindent}

\SymbolLegendDeclareReading{coordinates-cartesian}{\(X\times Y\)}%
  {Alle geordneten Paare \((x,y)\) mit \(x\in X\), \(y\in Y\).}
\SymbolLegendDeclareReading{coordinates-left}{\(Ae\)}%
  {Alle Werte \(x\star e\) mit \(x\in A\): mögliche erste Koordinaten.}
\SymbolLegendDeclareReading{coordinates-right}{\(eA\)}%
  {Alle Werte \(e\star x\) mit \(x\in A\): mögliche zweite Koordinaten.}
\SymbolLegendDeclareReading{coordinates-carrier}{\(\mathcal K_e\)}%
  {Der Koordinatenträger \(Ae\times eA\): alle Kombinationen.}
\SymbolLegendDeclareReading{coordinates-map}{\(\chi_e\)}%
  {Schickt \(x\) auf das Paar \((x\star e,e\star x)\).}
\SymbolLegendDeclareReading{coordinates-multiply}{\(\mu_e\)}%
  {Multipliziert ein Koordinatenpaar: \((p,q)\mapsto p\star q\).}
\SymbolLegendDeclareReading{coordinates-product}{\(\square_e\)}%
  {Paarprodukt: \((p,q)\square_e(r,s)=(p,s)\).}
\SymbolLegendDeclareID[prior]{coordinates-star}{\(\star\)}%
  {Gegebene assoziative Operation auf der Halbgruppe.}{SemigroupDef}
\SymbolLegendDeclareID[prior]{coordinates-composition}{\(g\circ f\)}%
  {Zuerst \(f\), dann \(g\): \((g\circ f)(x)=g(f(x))\).}{CompositionDef}
\SymbolLegendDeclareID[prior]{coordinates-bijection}{\(f\colon X\bij Y\)}%
  {Jedes Element von \(Y\) hat genau ein Urbild in \(X\).}{Bijektive Funktion}
\SymbolLegendDeclareID[prior]{coordinates-identity}{\(\Id_X\)}%
  {Identität auf \(X\): lässt jedes \(x\in X\) unverändert.}{IdentityFunctionDef}

\chapter*{Eine Halbgruppe als Rechteck lesen}
\hypertarget{coordinates.reading}{}

In einer Tabelle mit mehreren Zeilen und Spalten bestimmen
Zeile und Spalte gemeinsam ein Feld.
Jede der beiden Angaben allein lässt mehrere Felder offen; zusammen
legen sie genau eines fest. Wir möchten die Elemente einer Halbgruppe
ähnlich beschreiben. Ihre beiden Angaben sollen sich aber aus der
gegebenen Multiplikation gewinnen lassen, auch wenn die Elemente
zunächst gar keine Zeilen- und Spaltennummern besitzen.

Dazu wählen wir ein festes Element \(e\) als Bezugspunkt.
Für ein beliebiges \(x\) bilden wir zwei Produkte: einmal steht
\(e\) rechts, einmal links. Die
\emph{\(e\)-Koordinatenabbildung} ordnet \(x\) das Paar
\((x\CRStar e,e\CRStar x)\) zu. Der Buchstabe \(e\) erinnert
hier an das gewählte Bezugselement; wir verlangen kein Einselement
und keine neutrale Wirkung. Wie dieses Element eine Bezugszeile und
eine Bezugsspalte festlegen kann, zeigt gleich unser erstes Beispiel.

Zunächst ist das Paar nur ein Versuch, \(x\) durch zwei Angaben
zu erfassen. Entscheidend sind drei Fragen: Lassen sich verschiedene
Elemente unterscheiden? Kommt jede Kombination der möglichen Angaben
vor? Und lässt sich das Produkt zweier Elemente aus ihren Angaben
berechnen? Wir werden zeigen, dass zwei Gesetze dafür ausreichen.
Für alle \(x,y,z\) der Halbgruppe sollen gelten:
\[
  x\CRStar x=x,
  \qquad (x\CRStar y)\CRStar z=x\CRStar z.
\]
Das erste heißt \emph{Idempotenz}: Jedes Element bleibt beim
Multiplizieren mit sich selbst unverändert. Im zweiten, dem
\emph{Dreitermgesetz}, darf der mittlere Faktor aus jedem Produkt
dreier Faktoren gestrichen werden. Das Ergebnis ist genau das
Produkt des ersten mit dem letzten Faktor. Diese Gesetze gelten
nicht in jeder Halbgruppe; wir prüfen später auch Gegenbeispiele.

\section*{Ein Rechteck mit sechs Feldern}
\hypertarget{coordinates.reading.rectangle}{}

Wir beginnen mit zwei Zeilen und drei Spalten:
\[
  I=\{1,2\},\qquad J=\{a,b,c\},\qquad
  A=I\CRBin*{coordinates-cartesian}{\times}J.
\]
Das \emph{kartesische Produkt} enthält alle geordneten Paare, deren
erster Eintrag aus \(I\), deren zweiter aus \(J\) stammt.
Ein Paar \((i,j)\) ist hier das Feld in Zeile \(i\), Spalte \(j\).
Die Reihenfolge der Einträge gehört zur Information: Der erste
Eintrag bezeichnet immer die Zeile.
\[
  \begin{array}{c|ccc}
    &a&b&c\\\hline
    1&(1,a)&(1,b)&(1,c)\\
    2&(2,a)&(2,b)&(2,c)
  \end{array}
\]
Wir multiplizieren zwei Felder, indem wir die Zeile des linken und
die Spalte des rechten Feldes nehmen:
\[
  (i,j)\CRStar(k,\ell):=(i,\ell).
\]
Beispielsweise ist
\[
  (2,a)\CRStar(1,c)=(2,c),\qquad
  (1,c)\CRStar(2,a)=(1,a).
\]
Die Operation ist also nicht kommutativ. Wohl aber ist sie
assoziativ. Für drei beliebige Felder gilt
\[
 \begin{aligned}
  ((i,j)\CRStar(k,\ell))\CRStar(m,n)&=(i,\ell)\CRStar(m,n)=(i,n),\\
  (i,j)\CRStar((k,\ell)\CRStar(m,n))&=(i,j)\CRStar(k,n)=(i,n).
 \end{aligned}
\]
Beide Klammerungen liefern dasselbe Feld. Außerdem bleibt jedes
Produkt in \(A\). Damit ist dieses Rechteck eine Halbgruppe.

Wählen wir nun das Feld \(e=(1,a)\) als Bezugspunkt. Für
\(x=(2,c)\) erhalten wir
\[
  x\CRStar e=(2,a),\qquad e\CRStar x=(1,c).
\]
Die erste Rechnung bringt uns in die Bezugsspalte \(a\), behält aber
die Zeile von \(x\). Die zweite bringt uns in die Bezugszeile~\(1\),
behält aber die Spalte. Aus den beiden gefundenen Feldern lässt
sich \(x\) wiederherstellen:
\[
  (2,a)\CRStar(1,c)=(2,c)=x.
\]
Die neuen Koordinaten sind hier also selbst \emph{Elemente der
Halbgruppe}. Sie sind nicht einfach die ursprünglichen Einträge
\(2\) und \(c\), sondern die beiden Felder \((2,a)\) und
\((1,c)\). Zusammen enthalten sie dieselbe Information.

Das Basiselement \(e\) ist dabei kein neutrales Element. Schon
\((2,c)\CRStar e=(2,a)\ne(2,c)\) zeigt den Unterschied.
Es legt die Bezugszeile und die Bezugsspalte fest. Genauer: Ist
\(e=(i_0,j_0)\), dann gilt für jedes \(x=(i,j)\)
\[
  x\CRStar e=(i,j_0),\qquad e\CRStar x=(i_0,j).
\]
Die beiden Werte liegen also in der festen Spalte \(j_0\) und
der festen Zeile \(i_0\). Darum nennen wir \(e\) ein
\emph{Basiselement}: Es bestimmt die Bezugslinien, von denen aus
wir die Elemente beschreiben. Gemeint ist keine Vektorraumbasis.
Man könnte dasselbe gewählte Element auch \(b\) nennen und von
\(b\)-Koordinaten sprechen. Der Name und die Buchstabenwahl sind
hier Teil der Notation des Skripts; ihr mathematischer Hintergrund
ist die am Schluss beschriebene Rechteckdarstellung.

\section*{Welche Koordinaten gibt es in einer beliebigen Halbgruppe?}
\hypertarget{coordinates.reading.construction}{}

Jetzt sei \((A,\star)\) eine beliebige Halbgruppe und \(e\in A\).
Die Halbgruppenvoraussetzung aus
\FormulaRefAuto{SemigroupDef}[def] bedeutet: Jedes Produkt zweier
Elemente von \(A\) ist eindeutig bestimmt und liegt wieder in \(A\);
außerdem gilt das Assoziativgesetz. Die Wahl von \(e\) stellt hier
insbesondere sicher, dass \(A\) nicht leer ist. Idempotenz oder das
Dreitermgesetz nehmen wir zunächst noch nicht an.

Die beiden Mengen der möglichen Koordinaten heißen kurz
\[
 \begin{aligned}
  \CROrd*{coordinates-left}{Ae}&:=\{x\CRStar e\mid x\in A\},\\
  \CROrd*{coordinates-right}{eA}&:=\{e\CRStar x\mid x\in A\}.
 \end{aligned}
\]
Diese Schreibweise bezeichnet Mengen von Produktwerten.
\(Ae\) ist beispielsweise die Menge aller Ergebnisse, die entstehen,
wenn rechts immer dasselbe \(e\) steht. Beide Mengen liegen in \(A\),
weil die Operation abgeschlossen ist. Sie sind nicht leer: Der Wert
\(e\CRStar e\) gehört zu beiden.

\CRStatement{SemigroupCoordinateCarrierDef}{def}{Koordinatenträger}
Der \emph{Koordinatenträger} enthält jede Kombination einer
möglichen ersten mit einer möglichen zweiten Koordinate:
\[
  \CROrd*{coordinates-carrier}{\mathcal K_e}
    :=\CRLeft\CRTimes\CRRight.
\]
Dies ist die Schreibweise der vorstehenden Definition im Hauptband.
Die beiden Koordinaten werden hier unabhängig gewählt. Die Menge enthält also
auch Paare, deren erste und zweite Koordinate zunächst von
verschiedenen Ausgangselementen stammen.

\CRStatement{SemigroupCoordinateCarrierMembership}{thm}{Elementkriterium}
Ein Paar \((p,q)\) gehört genau dann zu \(\CRK\), wenn
\(p,q\in A\) sind und Elemente \(a,b\in A\) existieren mit
\[
  p=a\CRStar e,\qquad q=e\CRStar b.
\]
Zum Beweis lesen wir die Definition in beiden Richtungen.
Aus \((p,q)\in\CRLeft\CRTimes\CRRight\) folgt
\(p\in\CRLeft\), \(q\in\CRRight\); die Definitionen der
beiden Mengen liefern genau die genannten Zeugen \(a,b\).
Umgekehrt zeigen solche Zeugen, dass \(p\) in der ersten und
\(q\) in der zweiten Menge liegt, also ihr Paar im Produkt.
Das beweist \FormulaRefAuto{SemigroupCoordinateCarrierMembership}[thm]
(\CRProof{SemigroupCoordinateCarrierMembership}).

\CRStatement{SemigroupCoordinateCarrierNonempty}{thm}{Nichtleerer Koordinatenträger}
Zu jeder Halbgruppe mit gewähltem \(e\in A\) gehört mindestens
ein Koordinatenpaar. Zum Beweis liefern \(a=b=e\) das Paar
\[
  (e\CRStar e,e\CRStar e)\in\CRK.
\]
Damit ist auch der Koordinatenträger nicht leer. Dass
\(e\CRStar e=e\) wäre, wurde hierfür nicht gebraucht.
Dies ist \FormulaRefAuto{SemigroupCoordinateCarrierNonempty}[thm]
(\CRProof{SemigroupCoordinateCarrierNonempty}).

\section*{Das neue Produkt ist immer ein Rechteckprodukt}
\hypertarget{coordinates.reading.product}{}

\CRStatement{SemigroupCoordinateProductDef}{def}{Koordinatenprodukt}
Das Produkt zweier Koordinatenpaare übernimmt die erste Koordinate
des linken und die zweite Koordinate des rechten Paars:
\[
  (p,q)\CRBin*{coordinates-product}{\square_e}(r,s):=(p,s).
\]
Das Zeichen \(\CRSquare\) unterscheidet diese
neue Operation von \(\CRStar\), der gegebenen Multiplikation in
\(A\). Es handelt sich insbesondere nicht um das gewöhnliche
komponentenweise Produkt \((p\star r,q\star s)\).

\medskip\noindent\textbf{Warum ein Rechteck?}
Der entscheidende Punkt ist die freie Kombination: Zu zwei möglichen
Zeilen \(p,r\) und zwei möglichen Spalten \(q,s\) enthält
\(\CRK=\CRLeft\CRTimes\CRRight\) alle vier Ecken.
Für \(u=(p,q)\) und \(v=(r,s)\) sieht das so aus:
\[
 \begin{array}{c|cc}
     &q&s\\\hline
  p&u=(p,q)&u\CRSquare v=(p,s)\\
  r&v\CRSquare u=(r,q)&v=(r,s)
 \end{array}
\]
Das Produkt wählt die Ecke in der Zeile von \(u\) und der
Spalte von \(v\). Die umgekehrte Reihenfolge wählt die andere
gemischte Ecke. Fallen Zeilen oder Spalten zusammen, kann eine
solche Ecke auch mit einem Ausgangspunkt übereinstimmen.
Längen und Winkel spielen bei dieser Rechteckvorstellung keine Rolle.

Die bloße Beschreibung durch Paare würde dafür nicht genügen.
Bei den drei Paaren \((0,0),(0,1),(1,0)\) fehlt zum Beispiel
die vierte Ecke \((1,1)\). Die Regel würde aus \((1,0)\) und
\((0,1)\) gerade diese fehlende Ecke erzeugen und wäre auf der
Dreipunktmenge keine abgeschlossene Operation. Das kartesische
Produkt sichert, dass alle benötigten Ecken tatsächlich vorhanden sind.

\CRStatement{SemigroupCoordinateProductEvaluation}{thm}{Auswertung des Koordinatenprodukts}
Für zwei Paare \((p,q),(r,s)\) im Koordinatenträger ist
ihr Produkt genau das Paar \((p,s)\).

Warum ist die Regel auf ganz \(\CRK\) ausführbar? Liegen
\((p,q),(r,s)\in\CRK\), so ist \(p\in\CRLeft\) und
\(s\in\CRRight\). Folglich liegt auch \((p,s)\) im
Koordinatenträger. Ein geordnetes Paar hat eindeutig bestimmte
Einträge; das Ergebnis hängt daher nicht von einer Darstellung ab.
Die Regel definiert eine binäre Operation. Ihre Auswertung ist
gerade \FormulaRefAuto{SemigroupCoordinateProductEvaluation}[thm]
aus \FormulaRefAuto{SemigroupCoordinateProductDef}[def]
(\CRProof{SemigroupCoordinateProductEvaluation}).

\CRStatement{SemigroupCoordinateProductSemigroup}{thm}{Die Koordinaten bilden eine Halbgruppe}
Für jede Halbgruppe mit \(e\in A\) ist das Koordinatenprodukt
auf \(\CRK\) abgeschlossen und assoziativ.
Die Abgeschlossenheit haben wir gerade gezeigt; für die
Assoziativität vergleichen wir die beiden Klammerungen.
Jedes Element von \(\CRK=\CRLeft\CRTimes\CRRight\) ist ein
geordnetes Paar. Drei beliebige Elemente können wir daher als
\((p,q),(r,s),(t,u)\) schreiben. Die Zwischenprodukte
\((p,s)\) und \((r,u)\) gehören ebenfalls zu \(\CRK\), weil
ihre erste Koordinate in \(\CRLeft\), ihre zweite in
\(\CRRight\) liegt. Damit sind beide folgenden Rechnungen zulässig:
\[
 \begin{aligned}
  ((p,q)\CRSquare(r,s))\CRSquare(t,u)
    &=(p,s)\CRSquare(t,u)=(p,u),\\
  (p,q)\CRSquare((r,s)\CRSquare(t,u))
    &=(p,q)\CRSquare(r,u)=(p,u).
 \end{aligned}
\]
Die beiden Ergebnisse stimmen überein. Zusammen mit der eben
gezeigten Abgeschlossenheit beweist dies, dass
\((\CRK,\square_e)\) eine Halbgruppe ist:
\FormulaRefAuto{SemigroupCoordinateProductSemigroup}[thm]
(\CRProof{SemigroupCoordinateProductSemigroup}).

\CRStatement{SemigroupCoordinateProductIdempotence}{thm}{Idempotenz}
Jedes Koordinatenpaar ist bezüglich \(\CRSquare\) idempotent.
Jedes \(u\in\CRK\) lässt sich als \((p,q)\) schreiben. Daher
\[
  u\CRSquare u=(p,q)\CRSquare(p,q)=(p,q)=u.
\]
Dies beweist \FormulaRefAuto{SemigroupCoordinateProductIdempotence}[thm]
(\CRProof{SemigroupCoordinateProductIdempotence}).

\CRStatement{SemigroupCoordinateProductThreeTermLaw}{thm}{Dreitermgesetz}
Bei jedem Produkt dreier Koordinatenpaare darf der mittlere Faktor
gestrichen werden.
Für \(u=(p,q)\), \(v=(r,s)\) und \(w=(t,z)\) erhalten wir
\[
  (u\CRSquare v)\CRSquare w
   =(p,s)\CRSquare(t,z)=(p,z)
   =u\CRSquare w.
\]
Der mittlere Faktor \(v\) liefert keinen der beiden Einträge des
Endergebnisses. Damit ist
\FormulaRefAuto{SemigroupCoordinateProductThreeTermLaw}[thm]
bewiesen (\CRProof{SemigroupCoordinateProductThreeTermLaw}).

\CRStatement{SemigroupCoordinateProductRectangularIdentity}{thm}{Rechtecksidentität}
Für beliebige Koordinatenpaare \(u,v\) gilt
\((u\CRSquare v)\CRSquare u=u\).
Setzen wir im Dreitermgesetz \(w=u\), so folgt mit der Idempotenz
\[
  (u\CRSquare v)\CRSquare u=u\CRSquare u=u.
\]
Das ist \FormulaRefAuto{SemigroupCoordinateProductRectangularIdentity}[thm]
(\CRProof{SemigroupCoordinateProductRectangularIdentity}).
Nach der Zwischenstation \(v\) nimmt das letzte \(u\) seine Spalte
wieder mit, während die Zeile bereits vom ersten \(u\) stammt.
So kehrt die Rechnung im eben gezeichneten Rechteck zur Ecke \(u\)
zurück: \((p,q)\) geht zunächst nach \((p,s)\) und dann wieder
nach \((p,q)\). Daher die Bezeichnung \emph{Rechtecksidentität}.
Eine Identität bedeutet hier eine Gleichung, die für alle eingesetzten
Elemente gilt.

Alle vier Eigenschaften gelten für den Koordinatenträger jeder
Halbgruppe mit gewähltem Basiselement. Sie sagen noch nichts darüber,
ob die ursprüngliche Halbgruppe ebenfalls ein Rechteck ist.
Diese Verbindung muss erst die Koordinatenabbildung herstellen.

\par\Needspace{10\baselineskip}
\section*{Zerlegen und Zusammensetzen: zwei wohldefinierte Abbildungen}
\hypertarget{coordinates.reading.maps}{}

\CRStatement{SemigroupCoordinateMapDef}{def}{Die \(e\)-Koordinatenabbildung}
Die \emph{Koordinatenabbildung} ordnet jedem Element die beiden
durch Multiplikation mit \(e\) gewonnenen Angaben zu:
\[
 \begin{aligned}
  \CROrd*{coordinates-map}{\chi_e}&\colon A\longrightarrow\CRK,\\
  \CRChi(x)&:=(x\CRStar e,e\CRStar x).
 \end{aligned}
\]
\CRStatement{SemigroupCoordinateMapFunction}{thm}{Die Koordinatenabbildung ist eine Funktion}
Für jede Halbgruppe mit \(e\in A\) ordnet \(\CRChi\) jedem
Element von \(A\) genau ein Paar in \(\CRK\) zu.
Zum Beweis liegt für jedes \(x\in A\) die erste Koordinate in \(\CRLeft\),
die zweite in \(\CRRight\); in beiden Fällen dient \(x\) selbst
als Zeuge. Die beiden Produkte sind eindeutig bestimmt. Deshalb
gehört jedem \(x\) genau ein Paar im angegebenen Zielbereich.
Die Definition enthält bereits den angegebenen Funktionstyp;
unsere Überlegung überprüft zusätzlich seinen Zielbereich
(\CRProof{SemigroupCoordinateMapFunction}).

\CRStatement{SemigroupCoordinateMapEvaluation}{thm}{Auswertung der Koordinatenabbildung}
Das Bild eines beliebigen \(x\in A\) ist das Paar
\((x\CRStar e,e\CRStar x)\). Dies folgt unmittelbar aus der
Definitionsvorschrift
(\CRProof{SemigroupCoordinateMapEvaluation}).

\CRStatement{SemigroupCoordinateMultiplicationMapDef}{def}{Multiplikationsabbildung}
Für die entgegengesetzte Richtung bietet sich die
\emph{Multiplikationsabbildung} an:
\[
 \begin{aligned}
  \CROrd*{coordinates-multiply}{\mu_e}&\colon\CRK\longrightarrow A,\\
  \CRMu(p,q)&:=p\CRStar q.
 \end{aligned}
\]
\CRStatement{SemigroupCoordinateMultiplicationMapFunction}{thm}{Die Multiplikationsabbildung ist eine Funktion}
Jedes Koordinatenpaar wird durch \(\CRMu\) auf genau ein Element
der Ausgangshalbgruppe abgebildet.
Bei jedem Paar \((p,q)\in\CRK\) liegen beide Einträge in \(A\).
Die gegebene Operation liefert daher genau einen Wert \(p\CRStar q\)
in \(A\). Weil die Einträge des Paars eindeutig sind, ist auch diese
Vorschrift wohldefiniert
(\CRProof{SemigroupCoordinateMultiplicationMapFunction}).
Auch hier enthält die Definition bereits den Funktionstyp; die Rechnung
erklärt, weshalb jedes Koordinatenpaar einen Wert in \(A\) liefert.
\CRStatement{SemigroupCoordinateMultiplicationMapEvaluation}{thm}{Auswertung der Multiplikationsabbildung}
Für jedes Koordinatenpaar \((p,q)\) ist das Bild unter
\(\CRMu\) genau das ursprüngliche Produkt \(p\CRStar q\).
Auch dies ist unmittelbar die Definitionsvorschrift
(\CRProof{SemigroupCoordinateMultiplicationMapEvaluation}).

\par\medskip\noindent\fbox{%
\begin{minipage}{\dimexpr\linewidth-2\fboxsep-2\fboxrule\relax}
\textbf{Beispiel: die sechs Felder mit \(e=(1,a)\).}\par\smallskip
Wir verwenden wieder das Rechteck mit zwei Zeilen und drei Spalten
und wählen ausdrücklich dasselbe Basiselement \(e=(1,a)\).
Damit gilt
\[
 \begin{aligned}
  \CRLeft&=\{(1,a),(2,a)\},\\
  \CRRight&=\{(1,a),(1,b),(1,c)\},\\
  \CRChi(i,j)&=((i,a),(1,j)),\\
  \CRMu((i,a),(1,j))&=(i,j).
 \end{aligned}
\]
Es gibt zwei mögliche erste und drei mögliche zweite Koordinaten,
also sechs Koordinatenpaare. Jedes beschreibt genau ein Feld.
Bei \(x=(2,b)\) und \(y=(1,c)\) lässt sich auch das Produkt
vollständig in diesen neuen Koordinaten verfolgen. Die ursprüngliche
Regel \((i,j)\CRStar(k,\ell)=(i,\ell)\) benötigt kein
Basiselement. Deshalb steht \(x\CRStar y=(2,c)\) bereits fest,
bevor wir \(e\) wählen. Für die Koordinaten dieses Produkts und
der beiden Faktoren verwenden wir dagegen \(e=(1,a)\):
\[
 \begin{aligned}
  x\CRStar y&=(2,c),\\
  \CRChi(x\CRStar y)&=((2,a),(1,c)),\\
  \CRChi(x)\CRSquare\CRChi(y)
    &=((2,a),(1,b))\CRSquare((1,a),(1,c))\\
    &=((2,a),(1,c)).
 \end{aligned}
\]
Erst multiplizieren und dann Koordinaten nehmen liefert somit
dasselbe Ergebnis wie erst Koordinaten nehmen und dann das
Rechteckprodukt ausführen.
Der Index \(e\) am Zeichen \(\CRSquare\) benennt den
Koordinatenträger, auf dem gerechnet wird. Die Auswahlregel
\emph{erste Koordinate von links, zweite von rechts} selbst verwendet
keine zusätzliche Multiplikation mit \(e\).
\end{minipage}}
\par\medskip

\section*{Wann ist die Koordinatenabbildung ein Isomorphismus?}
\hypertarget{coordinates.reading.isomorphism}{}

Eine Funktion heißt \emph{injektiv}, wenn gleiche Bilder nur von
gleichen Ausgangselementen stammen. Sie heißt \emph{surjektiv}, wenn
jedes Element ihres Zielbereichs als Bild vorkommt. Erfüllt sie
beides, so ist sie \emph{bijektiv}: Jedes Zielelement besitzt genau
ein Urbild. Die Schreibweise
\[
  \CRChi\colon A\CRBij\CRK
\]
drückt diese Eigenschaft aus; sie stammt aus
\FormulaRefAuto{Bijektive Funktion}[def]. Für unsere Konstruktion
bedeutet sie, dass die beiden Koordinaten kein Element verwechseln
und keine mögliche Kombination auslassen.

Eine Bijektion ist erst dann ein \emph{Halbgruppenisomorphismus},
wenn sie außerdem die Multiplikation erhält:
\[
  \CRChi(x\CRStar y)
     =\CRChi(x)\CRSquare\CRChi(y)
       \qquad(x,y\in A).
\]
Beide Wege der Rechnung im vorigen Beispiel müssen also für
\emph{alle} Faktoren übereinstimmen.

\CRStatement{SemigroupCoordinateIsoDef}{def}{\(e\)-Koordinatenisomorphismus}
Die Koordinatenabbildung \(\CRChi\) heißt ein
\(e\)-Koordinatenisomorphismus, wenn \((A,\star)\) eine Halbgruppe
ist, \(e\in A\) gilt, \(\CRChi\) bijektiv von \(A\) auf
\(\CRK\) abbildet und die vorstehende Produktgleichung für alle
\(x,y\in A\) erfüllt. Damit fasst die Definition genau vier
Voraussetzungen zusammen.
Die vier folgenden Zugriffsaxiome im Hauptband machen diese
Bestandteile einzeln verfügbar; sie verlangen keine zusätzlichen
Eigenschaften.

\CRStatement{SemigroupCoordinateIsoIsSemigroupIso}{thm}{Ein Halbgruppenisomorphismus}
Ist \(\CRChi\) ein \(e\)-Koordinatenisomorphismus, so ist
\(\CRChi\) ein Halbgruppenisomorphismus von \((A,\star)\) auf
\((\CRK,\CRSquare)\).

\medskip\noindent\emph{Beweis.}
Unter der genannten Voraussetzung sind Quelle und Ziel Halbgruppen:
die Quelle nach Annahme, das Ziel nach
\FormulaRefAuto{SemigroupCoordinateProductSemigroup}[thm]. Dort wurden
Abgeschlossenheit und Assoziativität des Rechteckprodukts bewiesen.
Die Koordinatenabbildung ist
eine Funktion zwischen ihren Trägermengen und erhält die Operation.
Damit ist sie nach \FormulaRefAuto{SemigroupHomDef}[def] ein
Homomorphismus. Weil sie zusätzlich bijektiv ist, erfüllt sie
\FormulaRefAuto{SemigroupIsoDef}[def]. Das beweist
\FormulaRefAuto{SemigroupCoordinateIsoIsSemigroupIso}[thm]
(\CRProof{SemigroupCoordinateIsoIsSemigroupIso}).
Der Ausdruck \emph{\(e\)-Koordinatenisomorphismus} benennt in
diesem Werk also einen ganz bestimmten Halbgruppenisomorphismus,
nämlich den durch \(x\mapsto(x\star e,e\star x)\) gegebenen.

\section*{Das Kriterium: Zwei Gesetze reichen aus}
\hypertarget{coordinates.reading.criterion}{}

\CRStatement{SemigroupCoordinateIsoCriterion}{thm}{Das \(e\)-Koordinatenisomorphiekriterium}
Sei \((A,\star)\) eine Halbgruppe, \(e\in A\), und für alle
\(x,y,z\in A\) gelte
\[
  x\CRStar x=x,
  \qquad (x\CRStar y)\CRStar z=x\CRStar z.
\]
Dann ist \(\CRChi\) ein \(e\)-Koordinatenisomorphismus.
Dies ist \FormulaRefAuto{SemigroupCoordinateIsoCriterion}[thm].

\medskip\noindent\emph{Beweis.}
Die bereits nachgewiesenen Funktionstypen
\[
  \CRChi\colon A\longrightarrow\CRK,
  \qquad \CRMu\colon\CRK\longrightarrow A
\]
gelten schon für jede Halbgruppe mit \(e\in A\). Wir dürfen die
beiden Funktionen also hintereinander anwenden. Erst jetzt verwenden
wir die zusätzlichen Gesetze: Zunächst gewinnen wir jedes \(x\in A\)
aus seinen Koordinaten zurück und beweisen damit die Injektivität.
Dann konstruieren wir zu einem beliebigen Koordinatenpaar ein Urbild
und erhalten die Surjektivität. Zuletzt prüfen wir die
Produktverträglichkeit. Das ist auch die Reihenfolge der Beweistabelle.
Bei jeder Rechnung liegen die Faktoren in \(A\); die Abgeschlossenheit
sichert dies auch für ihre Zwischenprodukte. So sind die jeweiligen
Anwendungen von Assoziativität, Idempotenz und Dreitermgesetz zulässig.

\medskip\noindent\textbf{Erst zerlegen, dann zusammensetzen.}
Für ein beliebiges \(x\in A\) liegt \(\CRChi(x)\) in
\(\CRK\), also im Definitionsbereich von \(\CRMu\).
Die beiden Auswertungsformeln ergeben
\[
 \begin{aligned}
  \CRMu(\CRChi(x))
    &=(x\CRStar e)\CRStar(e\CRStar x)\\
    &=\bigl((x\CRStar e)\CRStar e\bigr)\CRStar x\\
    &=(x\CRStar e)\CRStar x\\
    &=x\CRStar x\\
    &=x.
 \end{aligned}
\]
Die erste Gleichheit setzt beide Definitionen ein. Die zweite ist
die rückwärts gelesene Assoziativität für die drei Elemente
\(x\CRStar e,e,x\), die sämtlich in \(A\) liegen.
In der dritten wird das Dreitermgesetz auf \(x,e,e\) angewandt;
in der vierten auf \(x,e,x\). Die letzte Gleichheit ist Idempotenz.
Jedes Ausgangselement wird nach den beiden Schritten wiedergewonnen.

\medskip\noindent\textbf{Warum die Rückgewinnung Injektivität liefert.}
Seien \(x,y\in A\) mit \(\CRChi(x)=\CRChi(y)\).
Beide Bilder gehören zum Definitionsbereich von \(\CRMu\).
Wenden wir diese Funktion auf die gleichen Bilder an, erhalten wir
gleiche Werte. Die eben für jedes Element bewiesene Rückrechnung ergibt
\[
  x=\CRMu(\CRChi(x))=\CRMu(\CRChi(y))=y.
\]
Weil \(x\) und \(y\) beliebig waren, ist \(\CRChi\) injektiv.
Für diesen Schluss haben wir die Surjektivität noch nicht vorausgesetzt.

\medskip\noindent\textbf{Ein beliebiges Koordinatenpaar und seine Zeugen.}
Sei nun \(u\in\CRK\) beliebig. Weil
\(\CRK=\CRLeft\CRTimes\CRRight\), lässt sich \(u=(p,q)\)
mit \(p\in\CRLeft\) und \(q\in\CRRight\) schreiben.
Das Elementkriterium liefert \(p,q\in A\) sowie Zeugen
\(a,b\in A\) mit \(p=a\CRStar e\) und \(q=e\CRStar b\).
Zuerst zeigen wir
\[
 \begin{aligned}
  p\CRStar e&=(a\CRStar e)\CRStar e=a\CRStar e=p,\\
  e\CRStar q&=e\CRStar(e\CRStar b)
              =(e\CRStar e)\CRStar b=e\CRStar b=q.
 \end{aligned}
\]
In der ersten Zeile wurde das Dreitermgesetz benutzt.
In der zweiten klammern wir zunächst um und verwenden dann
\(e\CRStar e=e\), die Idempotenz des Basiselements.
Die erste Schlussgleichheit \(p\CRStar e=p\) enthält den
Hilfszeugen \(a\) nicht mehr, die zweite \(e\CRStar q=q\)
den Hilfszeugen \(b\) nicht mehr. Wir haben außer den genannten
Darstellungen nichts über diese Zeugen angenommen. Die beiden
Gleichheiten gelten daher für die vorgegebenen Koordinaten,
unabhängig davon, welche passenden Zeugen wir gewählt haben.

\medskip\noindent\textbf{Das Produkt der Komponenten ist ein Urbild.}
Das sind genau die Vereinfachungen, die nach dem Zusammensetzen
benötigt werden. Weil \(p,q\in A\) sind, ist
\(p\CRStar q\in A\). Dieser Wert darf also in \(\CRChi\)
eingesetzt werden. Für die erste zurückgewonnene Koordinate gilt
\[
  (p\CRStar q)\CRStar e=p\CRStar e=p.
\]
Für die zweite gilt
\[
  e\CRStar(p\CRStar q)=(e\CRStar p)\CRStar q
                       =e\CRStar q=q.
\]
Die erste Rechnung benutzt das Dreitermgesetz für \(p,q,e\),
die zweite nach dem Umklammern für \(e,p,q\). Damit folgt
\[
 \begin{aligned}
  \CRChi(\CRMu(p,q))
    &=\bigl((p\CRStar q)\CRStar e,
                     e\CRStar(p\CRStar q)\bigr)\\
    &=(p,q).
 \end{aligned}
\]
Auch jedes beliebige Koordinatenpaar wird also wiedergewonnen.
Hier mussten die beiden Einträge nicht vom selben Ausgangselement
stammen: Gerade diese Freiheit ist für die Surjektivität entscheidend.

Für das ursprünglich beliebige \(u=(p,q)\) haben wir damit das
konkrete Urbild \(x=p\CRStar q\in A\) gefunden:
\(\CRChi(x)=u\). Die Hilfszeugen \(a,b\) werden für das Urbild
nicht mehr gebraucht. Jedes \(u\in\CRK\) besitzt folglich ein
Urbild; dies ist die Surjektivität. Zusammen mit der zuvor bewiesenen
Injektivität ergibt sie die Bijektivität.

In der Sprache der Funktionstheorie lauten die beiden Rückrechnungen
\[
  \CRMu\CRComp\CRChi=\CRId A,
  \qquad
  \CRChi\CRComp\CRMu=\CRId{\mathcal K_e}.
\]
Die erste Komposition ist eine Funktion von \(A\) nach \(A\),
die zweite von \(\CRK\) nach \(\CRK\). Die
\emph{Komposition} \(g\CRComp f\) bedeutet: zuerst \(f\)
anwenden, dann \(g\). Die \emph{Identität} \(\CRId X\) lässt
jedes Element von \(X\) unverändert. Weil die Funktionen auf jeder
Eingabe übereinstimmen, dürfen die punktweisen Gleichheiten als
Funktionsgleichheiten geschrieben werden; dies ist
\FormulaRefAuto{FunctionExtensionality}[thm].
Die allgemeine Regel
\FormulaRefAuto{MutuallyInverseCompanionBijection}[thm] fasst den
soeben ausgeführten Bijektivitätsbeweis zusammen. Insbesondere ist
\(\CRMu\) jetzt wirklich die Umkehrfunktion von \(\CRChi\).

\medskip\noindent\textbf{Verträglichkeit mit der Multiplikation.}
Seien schließlich \(x,y\in A\) beliebig. Wegen der Abgeschlossenheit
liegt \(x\CRStar y\) in \(A\); zugleich gehören
\(\CRChi(x)\) und \(\CRChi(y)\) zu \(\CRK\).
Damit sind sowohl das Bild des Produkts als auch das Rechteckprodukt
der beiden Bilder zulässig. Wir rechnen
\[
 \begin{aligned}
  \CRChi(x\CRStar y)
    &=\bigl((x\CRStar y)\CRStar e,
                     e\CRStar(x\CRStar y)\bigr)\\
    &=\bigl(x\CRStar e,(e\CRStar x)\CRStar y\bigr)\\
    &=\bigl(x\CRStar e,e\CRStar y\bigr)\\
    &=\bigl(x\CRStar e,e\CRStar x\bigr)
          \CRSquare\bigl(y\CRStar e,e\CRStar y\bigr)\\
    &=\CRChi(x)\CRSquare\CRChi(y).
 \end{aligned}
\]
Der Übergang zur zweiten Zeile benutzt in der ersten Koordinate das
Dreitermgesetz und in der zweiten Assoziativität. In der dritten
Zeile wird dort ebenfalls das Dreitermgesetz angewandt. Die vierte
Zeile ist die Definition des Rechteckprodukts.

Die Quelle ist nach Voraussetzung eine Halbgruppe, \(e\) gehört
zu ihr, und Bijektivität sowie Produktverträglichkeit sind bewiesen.
Die letzte Rechnung galt für jedes \(x\in A\) und jedes \(y\in A\).
Damit sind alle vier Bestandteile des Begriffs erfüllt. Das beweist
das Kriterium vollständig; seine formale Ableitung steht in der
\CRProof{SemigroupCoordinateIsoCriterion}.

\section*{Die Umkehrungen: Was ein Koordinatenisomorphismus erzwingt}
\hypertarget{coordinates.reading.converse}{}

Jetzt setzen wir umgekehrt voraus, dass ein
\(e\)-Koordinatenisomorphismus vorliegt. Seine Definition liefert
eine Halbgruppe \((A,\star)\), ein Element \(e\in A\), die
Bijektivität von \(\CRChi\) und die Operationserhaltung.
Jede Bijektion ist injektiv, wie
\FormulaRefAuto{F\colon A\bij B \vdash F\colon A\inj B}[thm]
festhält. Die bereits bekannten Gesetze im
Koordinatenträger können deshalb auf \(A\) zurückübertragen werden.
Dabei braucht man nicht zu raten, wie \(\CRMu\) aussieht.

\CRStatement{SemigroupCoordinateIsoIdempotence}{thm}{Idempotenz in der Ausgangshalbgruppe}
Ist \(\CRChi\) ein \(e\)-Koordinatenisomorphismus, so ist jedes
Element von \(A\) idempotent: Für alle \(x\in A\) gilt
\(x\CRStar x=x\).

\medskip\noindent\emph{Beweis.}
Sei \(x\in A\). Sein Bild liegt in \(\CRK\), wo jedes Element
idempotent ist. Produktverträglichkeit und dieses Zielgesetz ergeben
\[
  \CRChi(x\CRStar x)
   =\CRChi(x)\CRSquare\CRChi(x)=\CRChi(x).
\]
Sowohl \(x\CRStar x\) als auch \(x\) gehören zu \(A\).
Die Injektivität erzwingt daher \(x\CRStar x=x\).
Weil \(x\) beliebig war, gilt die Idempotenz auf ganz \(A\).
Das ist \FormulaRefAuto{SemigroupCoordinateIsoIdempotence}[thm]
(\CRProof{SemigroupCoordinateIsoIdempotence}).

\CRStatement{SemigroupCoordinateIsoThreeTermLaw}{thm}{Dreitermgesetz in der Ausgangshalbgruppe}
Ist \(\CRChi\) ein \(e\)-Koordinatenisomorphismus, so gilt in
\(A\) das Dreitermgesetz: Für alle \(x,y,z\in A\) ist
\((x\CRStar y)\CRStar z=x\CRStar z\).

\medskip\noindent\emph{Beweis.}
Seien \(x,y,z\in A\). Durch Abgeschlossenheit gehören
\(x\CRStar y\), \((x\CRStar y)\CRStar z\) und
\(x\CRStar z\) zu \(A\); außerdem liegen alle drei Bilder
\(\CRChi(x),\CRChi(y),\CRChi(z)\) in \(\CRK\).
Zweimalige Produktverträglichkeit liefert
\[
 \begin{aligned}
  \CRChi\bigl((x\CRStar y)\CRStar z\bigr)
    &=\CRChi(x\CRStar y)\CRSquare\CRChi(z)\\
    &=\bigl(\CRChi(x)\CRSquare\CRChi(y)\bigr)
                                      \CRSquare\CRChi(z)\\
    &=\CRChi(x)\CRSquare\CRChi(z)\\
    &=\CRChi(x\CRStar z).
 \end{aligned}
\]
In der dritten Zeile verwenden wir das Dreitermgesetz im
Koordinatenträger. Die letzte Gleichheit ist die Produktverträglichkeit
in umgekehrter Leserichtung. Beide Argumente der gleichen Bildwerte
liegen in \(A\); durch Injektivität folgt
\[
  (x\CRStar y)\CRStar z=x\CRStar z.
\]
Die drei Elemente waren beliebig. Damit ist
\FormulaRefAuto{SemigroupCoordinateIsoThreeTermLaw}[thm]
bewiesen (\CRProof{SemigroupCoordinateIsoThreeTermLaw}).

\CRStatement{SemigroupCoordinateIsoRectangularIdentity}{thm}{Rechtecksidentität in der Ausgangshalbgruppe}
Ist \(\CRChi\) ein \(e\)-Koordinatenisomorphismus, so gilt in
\(A\) die Rechtecksidentität: Für alle \(x,y\in A\) ist
\((x\CRStar y)\CRStar x=x\).

\medskip\noindent\emph{Beweis.}
Setzen wir im gerade gewonnenen Gesetz \(z=x\), so folgt für
alle \(x,y\in A\)
\[
  (x\CRStar y)\CRStar x=x\CRStar x=x.
\]
Die letzte Gleichheit benutzt die ebenfalls schon bewiesene
Idempotenz. Das ist
\FormulaRefAuto{SemigroupCoordinateIsoRectangularIdentity}[thm]
(\CRProof{SemigroupCoordinateIsoRectangularIdentity}).

Kriterium und Umkehrungen geben somit eine genaue Charakterisierung:
Für eine Halbgruppe und ein Element \(e\in A\) ist die
Koordinatenabbildung genau dann ein Isomorphismus, wenn Idempotenz
und Dreitermgesetz gelten. In üblicher Terminologie heißt eine
Halbgruppe mit diesen beiden Eigenschaften ein
\emph{Rechtecksband}. Das Wort \glqq Band\grqq{} bezeichnet hier
eine idempotente Halbgruppe, keinen Buchband.

\section*{Was von der Wahl des Basiselements abhängt}
\hypertarget{coordinates.reading.base}{}

Die beiden Gesetze des Kriteriums enthalten \(e\) gar nicht.
Liegt deshalb für ein Basiselement eine Koordinatenisomorphie vor,
so liefern die Umkehrungen diese Gesetze. Das Kriterium gilt dann
erneut für jedes andere \(f\in A\). Das Rechteck hängt als
algebraische Struktur nicht an einem ausgezeichneten Bezugspunkt;
die konkreten Koordinatenmengen ändern sich mit der Wahl.

Im Sechsfelderbeispiel wählen wir statt \(e=(1,a)\) nun
\(f=(2,b)\). Dann gilt für jedes \(x=(i,j)\)
\[
  x\CRStar f=(i,b),\qquad f\CRStar x=(2,j).
\]
Die neue Darstellung von \((2,c)\) lautet also
\(((2,b),(2,c))\), die alte lautete \(((2,a),(1,c))\).
Beide werden durch Multiplikation ihrer Einträge zu \((2,c)\).

Die Umrechnung lässt sich auch ohne die ursprünglichen Zeilen- und
Spaltennamen formulieren. Ein altes Koordinatenpaar \((p,q)\)
stellt das Element \(p\CRStar q\) dar. Seine neuen Koordinaten
bezüglich \(f\) sind daher
\[
 \begin{aligned}
  \chi_f(\CRMu(p,q))
    &=\bigl((p\CRStar q)\CRStar f,
                         f\CRStar(p\CRStar q)\bigr)\\
    &=\bigl(p\CRStar f,f\CRStar q\bigr).
 \end{aligned}
\]
Die erste Vereinfachung ist das Dreitermgesetz. In der zweiten
Koordinate klammern wir zuvor um. Die Rückumrechnung entsteht
durch Vertauschen der Rollen von \(e\) und \(f\), weil beide
Koordinatensysteme bijektiv dasselbe \(A\) darstellen.

\medskip\noindent\textbf{Zeilen und Spalten dürfen entarten.}
Es ist nicht erforderlich, dass es mindestens zwei von jeder Art gibt.
Bei der Operation \(x\CRStar y=x\) auf einer nichtleeren Menge
ist \(\CRLeft=A\), \(\CRRight=\{e\}\), und
\(\CRChi(x)=(x,e)\). Das Rechteck hat nur eine Spalte.
Die Rückrechnung lautet \(\CRMu(x,e)=x\CRStar e=x\).
Bei \(x\CRStar y=y\) ist entsprechend
\(\CRLeft=\{e\}\), \(\CRRight=A\),
\(\CRChi(x)=(e,x)\) und \(\CRMu(e,x)=e\CRStar x=x\).
Hier gibt es nur eine Zeile. Beide Operationen sind assoziativ,
idempotent und erfüllen das Dreitermgesetz: Jede Klammerung liefert
im ersten Fall den ersten, im zweiten den letzten Faktor.

\section*{Warum eine allgemeine Halbgruppe nicht genügt}
\hypertarget{coordinates.reading.counterexamples}{}

Die Formeln für \(\CRChi\) und \(\CRMu\) sind immer
wohldefiniert. Dass sie sich umkehren und das Produkt erhalten,
ist dagegen eine echte Aussage. Die folgenden Beispiele trennen
diese Fragen.

\medskip\noindent\textbf{Addition: zwei gleiche Angaben sind keine freien Koordinaten.}
Auf \(A=\mathbb N\) sei \(\star\) die Addition und \(e=0\).
Dann ist \(\CRLeft=\mathbb N=\CRRight\) und
\[
  \CRK=\mathbb N\CRTimes\mathbb N,\qquad
  \CRChi(n)=(n,n),\qquad \CRMu(p,q)=p+q.
\]
Die Abbildung ist injektiv, denn aus \((m,m)=(n,n)\) folgt
\(m=n\). Sie ist nicht surjektiv: Das Paar \((1,2)\) liegt
im Koordinatenträger, aber kein \((n,n)\) ist gleich \((1,2)\).
Außerdem scheitert die Produktverträglichkeit bereits bei \(1,2\):
\[
  \CRChi(1+2)=(3,3),\qquad
  \CRChi(1)\CRSquare\CRChi(2)=(1,1)\CRSquare(2,2)=(1,2).
\]
Auch die Rückrechnung misslingt: \(\CRMu(\CRChi(1))=2\ne1\).
Das neutrale Element der Addition ist somit kein Ersatz für die
Gesetze des Kriteriums.

\medskip\noindent\textbf{Minimum: Idempotenz allein reicht nicht.}
Sei \(A=\{0,1\}\) und \(x\star y=\min(x,y)\).
Das Produkt dreier Elemente ist bei beiden Klammerungen ihr
gemeinsames Minimum, also ist die Operation assoziativ.
Auch gilt \(\min(x,x)=x\). Mit \(e=1\) erhält man jedoch
\[
  \CRLeft=\CRRight=A,\qquad
  \CRChi(0)=(0,0),\quad\CRChi(1)=(1,1).
\]
Die beiden gemischten Paare \((0,1),(1,0)\) werden ausgelassen.
Das Dreitermgesetz ist tatsächlich falsch:
\[
  \min(\min(1,0),1)=0\ne1=\min(1,1).
\]
Wählen wir stattdessen \(e=0\), werden sogar beide Elemente auf
\((0,0)\) geschickt. Diese zweite Koordinatenabbildung ist surjektiv
auf den nun einelementigen Träger, aber nicht injektiv. Kein
Basiselement repariert das fehlende Dreitermgesetz.

\par\Needspace{16\baselineskip}
\medskip\noindent\textbf{Konstantes Produkt: Das Dreitermgesetz allein reicht nicht.}
Auf \(A=\{0,a\}\) mit \(a\ne0\) sei jedes Produkt gleich~\(0\).
Beide Klammerungen dreier Faktoren sind dann \(0\), und ebenso
ist das Produkt des ersten mit dem letzten Faktor \(0\).
Die Operation ist also assoziativ und erfüllt das Dreitermgesetz.
Aber \(a\CRStar a=0\ne a\): Die Idempotenz fehlt.
Für \(e=0\) gilt
\[
  \CRLeft=\{0\}=\CRRight,\qquad
  \CRK=\{(0,0)\},\qquad
  \CRChi(0)=\CRChi(a)=(0,0).
\]
Die Koordinatenabbildung erhält hier sogar die Operation und ist
surjektiv. Sie verliert dennoch die Unterscheidung zwischen \(0\)
und \(a\) und kann deshalb kein Isomorphismus sein.

Die drei Beispiele erklären die Aufgaben im Beweis: Die Rückrechnung
auf \(A\) verhindert das Zusammenfallen verschiedener Elemente;
die Rückrechnung auf \(\CRK\) stellt jede Koordinatenkombination
bereit; die Produktgleichung sichert schließlich, dass die neue
Darstellung dieselbe Multiplikation beschreibt.

\clearpage
\section*{Historischer Abriss: Von Halbgruppen zu Rechteckskoordinaten}
\hypertarget{coordinates.reading.history}{}

Die Strukturtheorie der Halbgruppen untersucht, welche Formen eine
assoziative Multiplikation ohne die zusätzlichen Forderungen einer
Gruppe annehmen kann. Einen frühen grundlegenden Beitrag lieferte
David Rees 1940 mit \emph{On semi-groups}. Er stellte die
Untersuchung der Struktur solcher Systeme ausdrücklich in den
Zusammenhang algebraischer Methoden.\footnote{D. Rees,
\emph{On semi-groups}, Proceedings of the Cambridge Philosophical
Society 36 (1940), S.~387--400.
\href{https://doi.org/10.1017/S0305004100017436}{Originalveröffentlichung beim Verlag}.}
Unser Rechteckmodell gehört zu diesem größeren Anliegen, konkrete
Koordinaten für abstrakte Multiplikationsgesetze zu finden.

1951 veröffentlichte J.~A. Green \emph{On the Structure of
Semigroups}. Die nach ihm benannten Relationen vergleichen Elemente
anhand der von ihnen erzeugten einseitigen Ideale. In einem
Rechtecksband lassen sich diese Beziehungen als gemeinsame Zeilen
und gemeinsame Spalten sehen.\footnote{J.~A. Green,
\emph{On the Structure of Semigroups}, Annals of Mathematics,
2.~Serie, 54 (1951), S.~163--172,
\href{https://doi.org/10.2307/1969317}{Originalveröffentlichung}.
Zur Zeilen- und Spaltendeutung siehe auch J.~E. Leech,
\emph{Noncommutative Lattices}, 2020, Abschnitt~1.2, S.~21--22.}
Die Koordinatenidee ist damit zugleich ein anschaulicher Einstieg in
eine allgemeinere Beschreibung der inneren Struktur von Halbgruppen.

In der heutigen Darstellung bezeichnet \emph{Rechtecksband} genau
die hier behandelte idempotente Struktur mit Dreitermgesetz.
Die Darstellung durch die Mengen \(\CRLeft\) und \(\CRRight\) ist ein
bekannter Struktursatz. Jonathan E. Leech formuliert ihn in
\emph{Noncommutative Lattices} als Satz~1.2.4: Die Multiplikation
der beiden Koordinaten liefert den Isomorphismus zum ursprünglichen
Rechtecksband.\footnote{Jonathan E. Leech,
\emph{Noncommutative Lattices: Skew Lattices, Skew Boolean Algebras
and Beyond}, University of Primorska Press, Koper 2020,
Abschnitt~1.2, Satz~1.2.4, S.~21--22.
\href{https://www.hippocampus.si/ISBN/978-961-293-027-1.pdf}{Frei zugängliche Verlagsausgabe},
\href{https://doi.org/10.26493/978-961-293-027-1}{DOI}.}
Auch Leech bezeichnet das beliebig gewählte Element mit \(e\).
Weil er ein Rechtecksband betrachtet, gilt \(e\CRStar e=e\);
das ist Idempotenz und fordert keine neutrale Wirkung auf andere
Elemente.
Die hier bewiesene Abbildung \(\CRChi\) ist die ausdrücklich
angegebene Umkehrung dieser Multiplikationsabbildung.

Der Name \emph{\(e\)-Koordinatenisomorphismus} und die gemeinsame
Notation \(\CRK,\CRSquare,\CRChi,\CRMu\) sind die Bezeichnungen
dieses Werks. Damit wird kein historisch eigenständiger Satz dieses
Namens einem Autor zugeschrieben. Der mathematische Gehalt ist die
Rechteckdarstellung: Zwei durch Multiplikation mit einem Basiselement
gewonnene Angaben bestimmen jedes Element eindeutig, und beim Produkt
bleiben die erste Angabe links und die zweite rechts erhalten.
